\chapter{Các service nghiệp vụ và module common}
\label{ch:services}

Các chương trước bàn về những service hạ tầng. Chương này vẽ bản đồ ba
service làm việc thật, cộng với module \code{common} dùng chung --- và
qua chúng, hai cách microservice nói chuyện với nhau.

\section{course-service: con ngựa thồ}

Khóa học, học viên, bài tập, ghi danh, bài nộp --- CRUD kinh điển trên
Postgres (Chương~7), cộng thêm file đính kèm lưu trong MinIO
(Chương~9). Nó \emph{phát} sự kiện \code{ts.assignment.created} và
\emph{lắng nghe} \code{ts.user.activated} --- mẫu đáng để ý là nó phản
ứng với sự kiện nghiệp vụ của service khác mà không bao giờ gọi service
đó.

\section{auth-service: danh tính}

Mọi thứ Chương~5 đã mô tả: đăng ký kèm kích hoạt qua e-mail, đăng nhập
bằng mật khẩu và Google, 2FA TOTP, phát hành JWT, refresh token. Nó phát
ba sự kiện: \code{ts.user.registered}, \code{ts.user.activated},
\code{ts.user.admin-provisioned}.

\section{dictionary-service: anh nhỏ}

Tra cứu từ vựng trên MongoDB (Chương~8). Đơn giản có chủ ý --- nó tồn
tại một phần để chứng minh kiến trúc chứa được một document store, một
phần làm service ít rủi ro để thử mọi thứ trước.

\section{notification-service: consumer thuần túy}

Không có REST API đáng kể, không cơ sở dữ liệu. Nó đăng ký ba topic và
biến sự kiện thành e-mail (Maildev khi dev, một SMTP relay thật về sau):

\begin{center}
\small
\begin{tabular}{@{}lll@{}}
\toprule
Topic & Producer & Mail gửi đi \\
\midrule
\code{ts.user.registered} & auth-service & link kích hoạt \\
\code{ts.user.admin-provisioned} & auth-service & thông báo tài khoản \\
\code{ts.assignment.created} & course-service & thông báo bài tập mới \\
\bottomrule
\end{tabular}
\end{center}

Vì đầu vào \emph{duy nhất} của nó là Kafka, đây là service duy nhất
không thể chạy hữu ích trước khi có broker --- đó là lý do manifest
Kubernetes của nó đi cùng cột mốc Kafka (Chương~25) chứ không cùng các
service khác.

\section{Đồng bộ và bất đồng bộ, trong codebase này}

Dự án dùng cả hai kiểu tương tác, mỗi kiểu đúng chỗ của nó:

\begin{description}
  \item[Đồng bộ (HTTP qua gateway)] khi bên gọi cần câu trả lời để đi
  tiếp: frontend lấy danh sách khóa học. Ràng buộc về thời gian --- cả
  hai bên phải đang chạy.
  \item[Bất đồng bộ (sự kiện qua Kafka)] khi bên gọi chỉ cần ghi nhận
  một sự thật: một đăng ký vừa xảy ra; ai quan tâm thì hành động.
  auth-service không biết notification-service tồn tại. Tách rời về thời
  gian --- nếu consumer sập, mail đi muộn chứ không mất.
\end{description}

Quy tắc ngón tay cái codebase tuân theo: \emph{lệnh} (command) cần câu
trả lời thì dùng HTTP; \emph{sự thật} mà người khác có thể phản ứng thì
là sự kiện, đặt tên ở thì quá khứ (\code{UserRegistered},
\code{AssignmentCreated}).

\section{Module common}

\begin{lstlisting}
com/ts/common/
  dto/        CourseDto, StudentDto, UserDto, ...
              UserRegisteredEvent, UserActivatedEvent,
              AssignmentCreatedEvent, AdminProvisionedUserEvent
  exception/  ApiException, ErrorResponse
  security/   JwtConstants     <- header names, claim names
\end{lstlisting}

Ba kiểu chia sẻ, với ba hồ sơ rủi ro khác nhau:

\begin{itemize}
  \item \textbf{Các record sự kiện} là \emph{hợp đồng} giữa producer và
  consumer. Chia sẻ class bảo đảm hai bên nhất trí --- với cái giá là
  phải triển khai đồng bộ khi hình dạng thay đổi. (Lựa chọn thay thế
  trong công nghiệp --- schema registry --- xuất hiện ở Chương~11.)
  \item \textbf{\code{JwtConstants}} chia sẻ \emph{tên} header và claim
  giữa gateway và các service. Một tên header gõ sai sẽ hỏng âm thầm
  --- request đơn giản là tới nơi dưới dạng vô danh --- nên gom về một
  chỗ là khoản bảo hiểm rẻ.
  \item \textbf{DTO và exception} là chia sẻ cho tiện: dễ chịu trong
  monorepo, và là thứ đầu tiên nên nhân-bản-thay-vì-chia-sẻ nếu các
  service có ngày tách ra thành các repository riêng, vì các class tiện
  ích dùng chung chính là cách những service ``độc lập'' rốt cuộc lại
  phải phát hành cùng nhau.
\end{itemize}

\begin{decision}[có nên có module dùng chung không?]
Những người theo chủ nghĩa thuần túy microservice cấm chia sẻ code
giữa các service. Lựa chọn thuần túy --- mỗi service tự định nghĩa bản
sao của mọi sự kiện --- lại cho phép trôi lệch đúng ở chỗ sự nhất trí là
bắt buộc. Dự án này vạch ranh giới ở: chia sẻ \emph{hợp đồng} (sự kiện,
tên header), nghĩ kỹ trước khi chia sẻ bất cứ thứ gì có hành vi.
\code{jib.skip=true} trong POM của common đánh dấu nó là module duy
nhất là thư viện, không phải thứ để triển khai.
\end{decision}

\begin{pitfall}[bẫy triển khai đồng bộ (lockstep)]
Triệu chứng: sau khi thêm một trường vào \code{UserRegisteredEvent} và
chỉ triển khai auth-service, notification-service bắt đầu ném lỗi giải
tuần tự trên mọi thông điệp --- và vì nó dùng chung consumer group, nó
cứ thử lại mãi cùng một thông điệp hỏng. Nguyên nhân: hợp đồng thay đổi
chỉ ở một phía, và topic giờ chứa sự kiện ở hai hình dạng. Các quy tắc
ngăn chuyện đó: chỉ \emph{thêm} trường tùy chọn, không bao giờ đổi tên
hay xóa; triển khai consumer trước producer; và cấu hình xử lý lỗi cho
consumer (Chương~11) để một thông điệp độc không làm nghẽn cả topic.
\end{pitfall}

\begin{quiz}
\begin{enumerate}
  \item Với mỗi trường hợp, nói HTTP hay sự kiện là phù hợp và vì sao:
  (a) frontend cần danh sách học viên; (b) một bài nộp vừa được chấm và
  học viên cần được báo; (c) course-service phải xác minh một JWT.
  \item Vì sao notification-service là service duy nhất có manifest
  Kubernetes phải chờ cột mốc Kafka, trong khi auth-service cũng dùng
  Kafka?
  \item Nêu một thứ thuộc về \code{common} và một thứ không bao giờ nên
  ở đó, kèm lập luận.
  \item Quy tắc thứ tự triển khai nào làm cho thay đổi schema sự kiện
  an toàn, và bản thân thay đổi đó phải có tính chất gì?
\end{enumerate}
\end{quiz}
